\lecture{4}{26 May. 14:18}{Fourth Lecture}

\subsubsection{Infix to Postfix Conversion Implementation}
Consider the relative precedence of operations according to \textbf{BODMAS}
rule. 

\begin{enumerate}
  \item \textbf{Initialize} an empty stack for operators, \(S\), and an empty output string (or list) for the postfix expression, \(P\).
  \item \textbf{Scan} the infix expression \(E\) from left to right, one symbol at a time.
  \item \textbf{For each symbol \(t\) in \(E\):}
    \begin{enumerate}
      \item If \(t\) is an \textbf{operand} (letter or digit), append \(t\) to \(P\).
      \item If \(t = (\), push \(t\) onto \(S\).
      \item If \(t = )\), pop and append all operators from \(S\) to \(P\) until a left parenthesis \((\) is encountered. Pop and discard the left parenthesis.
      \item If \(t\) is an \textbf{operator} (\(+\), \(-\), \(*\), \(/\), \(^\)):
        \begin{enumerate}
          \item While \(S\) is not empty, and the operator at the top of \(S\) has greater or equal precedence than \(t\) (and is not a left parenthesis), pop from \(S\) and append to \(P\).
          \item Push \(t\) onto \(S\).
        \end{enumerate}
    \end{enumerate}
  \item \textbf{After scanning} the entire infix expression, pop and append any remaining operators from \(S\) to \(P\).
  \item The output \(P\) is the postfix expression.
\end{enumerate}

\textbf{We define functions to handle precedence prehand and functions for handling expressions.}

\begin{listing}[H]
  \inputminted{python}{Code/expressions.py}
  \caption{Expressions Handling}
  \label{lst:postfix-eval-exp-handl}
\end{listing}

\begin{algorithm}[H]
  \caption{Infix to Postfix Conversion}
  \label{algo:infix-2-postfix-conversion}
  \KwIn{Infix expression \(E\)}
  \KwOut{Postfix expression \(P\)}

  Initialize an empty stack \(S\)\;
  Initialize an empty output string \(P\)\;
  \For{each token \(t\) in \(E\) (from left to right)}{
    \eIf{\(t\) is an operand}{
      Append \(t\) to \(P\)\;
    }{
      \eIf{\(t = (\)}{
        Push \(t\) onto \(S\)\;
      }{
        \eIf{\(t = )\)}{
          \While{top of \(S \neq (\)}{
            Pop from \(S\) and append to \(P\)\;
          }
          Pop \((\) from \(S\)\;
        }{
          \While{\(S\) is not empty \textbf{and} precedence(\(t\)) \(\leq\) precedence(top of \(S\)) \textbf{and} top of \(S \neq (\)}{
            Pop from \(S\) and append to \(P\)\;
          }
          Push \(t\) onto \(S\)\;
        }
      }
    }
  }
  \While{\(S\) is not empty}{
    Pop from \(S\) and append to \(P\)\;
  }
  \Return \(P\)\;
\end{algorithm}

\begin{listing}[H]
  \inputminted{python}{Code/infix_to_postfix.py}
  \caption{Infix to Postfix Conversion}
  \label{lst:infix-2-postfix-conversion}
\end{listing}


\subsection{Postfix Evaluation}

\begin{enumerate}
  \item Scan the expression from left to right until we encounter any operator.
  \item Perform the operation
  \item replace the expression with its computed value.
  \item Repeat the steps from \textbf{1} to \textbf{3} until no more
    operators exist.
\end{enumerate}

\begin{algorithm}
  \KwData{exp: Expression to evaluate}
  \KwResult{Result of evaluation of the expression}
  \Begin{
    stack \(\leftarrow\) empty stack\;
    \For{char \emph{in} exp}{
      \eIf{char is operator}{
        a \(\leftarrow\) stack.pop()\;
        b\(\leftarrow\)stack.pop()\;
        \(o \leftarrow \text{char}\)\;
        push a operated(\(o\)) on b to stack\;
      }{
        push char to stack\;
      }
    }
    \Return{pop stack}
  }
  \caption{Postfix Evaluation}
  \label{algo:postfix-eval}
\end{algorithm}

\begin{eg}
  Postfix Expression: \(34*25*+\)
  \begin{table}[H]
    \centering
    \begin{tabular}{|c|c|c|}
      \hline
      Input&Stack&Action\\
      \hline
      \(34*25*+\)&empty&Push 3\\
      \(4*25*+\)&3&Push 4\\
      \(*25*+\)&3,4& Pop 3, 4. Perform \(3*4\). Push 12\\
      \(25*+\)&12&Push(2)\\
      \(5*+\)&12, 2& Push 5\\
      \(*+\)&12, 2, 5&Pop 2, 5. Perform \(2*5\), Push 10\\
      \(+\)&12, 10&Pop 12, 10. Perform \(12+10\), Push 22\\
           &22&22 is the evaluation.\\
           \hline
    \end{tabular}
  \end{table}
\end{eg}

\subsubsection{Postfix Evaluation Implementation}
We have previously made a \href{sec:stack-impl}{Stack}. We have used the helper
functions made previously. Then the function for evaluating postfix expression is defined.

\begin{listing}[H]
  \inputminted{python}{Code/eval_postfix.py}
  \caption{Postfix Evaluation}
  \label{lst:postfix-eval}
\end{listing}
